\chapter{疼痛：警报信号}
\chaptermark{疼痛}
\label{sec:pain}

\section{是警报，不是测量}

第~\ref{sec:sensors}~章的所有传感器都在测量世界：光有多少，声音多高，压力多大。疼痛则不同。它报告的不是“那里有什么”，而是“危险，放下一切”。对工程师来说，这不是测量通道，而是\emph{警报信号}：一条优先级最高的线路，必须能穿透机器当前正在做的任何工作。

警报与测量通道有三点不同，而疼痛三点都具备。第一，它很难被适应压低：痛觉传感器的适应远弱于触觉传感器，轻度损伤后甚至变得更敏感~\cite{LaMotte1982hyper}；强烈加热后响应减弱，只在其中一类中测到过~\cite{Treede1995heat}。光传感器在恒定背景下会沉默（图~\ref{fig:adaptation}），而一分钟后就沉默的警报毫无用处。第二，它的阈值很高：痛觉传感器只在刺激接近危险时才发放，例如对加热，不同传感器的阈值从$41^\circ$到$46$--$48^\circ$不等，取决于类型~\cite{Treede1995heat, Weidner1999cfib}。第三——这也是本章的重点——警报的音量不仅由传感器决定，也由机器本身决定。同一处伤口在不同情境下疼得不一样。让我们看看它是由哪些已经熟悉的部件组成的。

痛觉传感器是\nt{GLUT}神经元，和第~\ref{sec:sensors}~章的其他感觉神经元一样。其中一部分在释放谷氨酸的同时还释放附录~\ref{sec:othermediators}中的P物质，它缓慢地增强传递。

\section{两根导线：快痛与慢痛}

撞一下手指，你会感到两次疼：先是短促尖锐的，一秒后是钝而持久的。这是两根不同的导线。快导线有绝缘，以$5$到$30$米每秒的速度$v$传导脉冲；慢导线细而裸露，速度为$0.5$--$2$米每秒。来自脚部的信号要走约一米，到达时间
\begin{empheq}[box=\keybox]{equation}
t \;=\; \frac{L}{v}
\label{eq:paintime}
\end{empheq}
对快导线是几十毫秒，对慢导线约一秒。两根导线传送两条消息。快导线报告\emph{哪里}和\emph{何时}：它的脉冲到得更早、更准，如同第~\ref{sec:code}~章的时间编码。慢导线报告\emph{有多糟}，它持久的钝痛使机器在损伤未消退时一直处于警报模式。

\section{脊髓中的闸门}

疼痛信号首先到达脊髓，并在那里通过一道\emph{闸门}~\cite{MelzackWall1965}；这道闸门的具体神经元是在相对晚近才找到的~\cite{Duan2014}。把疼痛继续传往大脑的输出\nt{GLUT}神经元上，汇聚着痛觉导线和触觉导线。旁边还有一个抑制性中间神经元，\nt{GABA}或\nt{GLYC}，它被触觉兴奋（图~\ref{fig:gate}）。触摸关闭闸门。在最初的闸门理论~\cite{MelzackWall1965}中，疼痛反过来关闭这个中间神经元，于是强烈的疼痛打开闸门；这是理论的假设，在已找到的闸门神经元上并未直接证实。

\begin{figure}[t]
\centering
\begin{tikzpicture}[>=Stealth,
  nrn/.style={circle,draw,very thick,minimum size=1.0cm,inner sep=0pt,font=\scriptsize},
  inh/.style={-{Bar[width=7pt]},thick,red!70!black},
  bc/.style={->,thick,densely dashed,orange!85!black}]
\node[nrn,fill=blue!6] (Z) at (6,0) {\nt{GLUT}};
\node[nrn,fill=red!10] (Q) at (3.6,-1.6) {\nt{GABA}};
\draw[->,very thick,red!70!black] (0,0.6) node[left,font=\small,black,align=right] {疼痛 $\nu$} -- (5.45,0.15);
\draw[->,very thick,blue!60!black] (0,-2.6) node[left,font=\small,black,align=right] {触摸 $\mu$} -- (2.9,-2.0);
\draw[->,thick,blue!60!black] (1.8,-2.35) -- (5.5,-0.35) node[pos=0.8,below right=-2pt,font=\scriptsize] {$w_{z\mu}$};
\draw[inh] (1.6,0.5) -- (3.35,-1.1) node[pos=0.5,left=1pt,font=\scriptsize] {$w_{q\nu}$};
\draw[inh] (Q) -- (Z) node[pos=0.5,above left=-1pt,font=\scriptsize] {$\beta$};
\draw[->,very thick] (Z) -- (8.2,0) node[right,font=\small,align=left] {送往大脑：$z$};
\draw[bc] (6,2.1) node[above,font=\scriptsize,black,align=center] {来自上方：\nt{SERO}、\nt{NORA}、阿片类} -- (Z.north) node[pos=0.45,right,font=\scriptsize,black] {增益 $g$};
\end{tikzpicture}
\caption{按模型~\eqref{eq:gate}的脊髓疼痛闸门。输出\nt{GLUT}神经元接收疼痛$\nu$和来自触摸$\mu$的微弱兴奋。抑制性中间神经元（\nt{GABA}或\nt{GLYC}）被触摸兴奋，并按闸门理论的假设被疼痛关闭。增益$g$经广播通道从上方到来。}
\label{fig:gate}
\end{figure}

像第~\ref{sec:gaba}~章那样，用频率把它写出来。中间神经元的频率$q$和输出频率$z$为
\begin{empheq}[box=\keybox]{equation}
q \;=\; \bigl[w_{q\mu}\,\mu + q_0 - w_{q\nu}\,\nu\bigr]_+ ,
\qquad
z \;=\; \bigl[\,g\,(\nu + w_{z\mu}\,\mu) - \beta\,q\,\bigr]_+ ,
\label{eq:gate}
\end{empheq}
其中$\nu$是疼痛输入，$\mu$是触摸输入，$w$是连接权重（第一个下标是接收方，第二个是发送方），$\beta$是抑制强度，$q_0$是中间神经元自身的背景活动，$g$是由上方设定的增益（见下文）。这就是第~\ref{sec:gaba}~章的否决~\eqref{eq:veto}：“有疼痛，但没有触摸”——只有一处修正，作为假设取自闸门理论：强烈的疼痛自己关闭否决神经元。

我们在$w_{q\mu}=1$、$q_0=0.2$、$w_{q\nu}=0.5$、$w_{z\mu}=0.3$、$\beta=1.5$时计算了模型（图~\ref{fig:gatecurves}）。没有触摸时，疼痛从$\nu\approx0.18$开始通过。有触摸时，阈值升到$0.86$，在$\nu=1$时输出降为四分之一，从$1$降到$0.25$。而在剧痛$\nu=2$时，触摸已经无济于事：闸门大开。生活中也是这样：揉一揉撞伤处有帮助，严重创伤时则不然。单纯的触摸，没有疼痛，不会通过闸门：触摸不会引发警报。

\begin{figure}[t]
\centering
\begin{tikzpicture}[>=Stealth,x=5.2cm,y=1.35cm]
\draw[->,gray!70!black] (0,0) -- (2.12,0) node[right,font=\small,black] {疼痛 $\nu$};
\draw[->,gray!70!black] (0,0) -- (0,4.3) node[left,font=\small,black] {$z$};
\foreach \t in {0,0.5,1,1.5,2}{\draw[gray!60!black] (\t,-0.05) -- (\t,0.05); \node[below,font=\scriptsize] at (\t,0) {\t};}
\foreach \f in {1,2,3,4}{\draw[gray!60!black] (-0.01,\f) -- (0.01,\f); \node[left,font=\scriptsize] at (0,\f) {\f};}
\draw[very thick,red!70!black] plot file {figures/pain_gate_notouch.dat};
\draw[very thick,blue!60!black] plot file {figures/pain_gate_touch.dat};
\draw[very thick,orange!85!black,densely dashed] plot file {figures/pain_gate_sens.dat};
\draw[very thick,red!70!black] (0.08,3.9) -- (0.2,3.9) node[right,font=\scriptsize,black] {无触摸};
\draw[very thick,blue!60!black] (0.08,3.45) -- (0.2,3.45) node[right,font=\scriptsize,black] {有触摸};
\draw[very thick,orange!85!black,densely dashed] (0.08,3.0) -- (0.2,3.0) node[right,font=\scriptsize,black] {敏化后，$g=2$};
\end{tikzpicture}
\caption{闸门~\eqref{eq:gate}的输出与疼痛强度的关系。触摸提高阈值，减弱轻度和中度疼痛，但对剧痛无效。敏化（虚线）使增益加倍：同样的疼痛以两倍的音量传出，阈值则降低。}
\label{fig:gatecurves}
\end{figure}

\section{来自上方的音量旋钮}

闸门不仅受下方控制。从脑干下行到闸门的通路会改变式~\eqref{eq:gate}中的增益$g$：\nt{SERO}、\nt{NORA}以及附录~\ref{sec:othermediators}中的阿片肽~\cite{Fields2004}。这正是第~\ref{sec:serocomputer}--\ref{sec:acet}~章的那些广播通道，只不过在这里它们的旋钮控制警报的音量。它们能把旋钮往两个方向拧：同样的下行回路既能减轻疼痛，也能加重疼痛。

机器为什么要调低警报？因为警报是一种中断，而中断并不总是合时宜。逃避捕食者的动物不应因为爪子受伤而停下：先跑，再舔伤口。对缓解的预期也会减轻疼痛，而阻断阿片受体就能消除这一效应的一部分~\cite{Levine1978}：在大脑中算出的预期，沿阿片通道下传到闸门。这幅图景本身已经测量；大脑究竟如何决定音量应该多大，还是假说。

\section{敏化：会学习的警报}

损伤之后，闸门的输出神经元变得更敏感：同样的输入产生更大的输出，阈值降低~\cite{Woolf1983}。在模型~\eqref{eq:gate}中，这是增益$g$的增大，最简单的描述是一个由疼痛本身充电的慢RC电路：
\begin{empheq}[box=\keybox]{equation}
\tau_s\,\frac{\dd g}{\dd t} \;=\; -(g-1) + k_s\,z ,
\label{eq:sensitization}
\end{empheq}
其中$\tau_s$是敏感性恢复正常所需的时间；它比疼痛本身持续得更久，因为在损伤性刺激停止之后敏化仍然持续~\cite{Woolf1983}；$k_s$是充电强度。只要组织受损，闸门的输出就使$g$保持在1以上，伤口附近发生的一切都被放大传出（图~\ref{fig:gatecurves}中的虚线：$g=2$时阈值降到$0.11$，输出加倍）。这是一种合理的警报设置：伤未愈合时，宁可多加小心。

在工程师看来，这是带有慢泄漏的正反馈：疼痛提高增益，增益又加重疼痛。只要环路较弱，伤口愈合后$g$就会恢复正常。但如果环路很强，警报可能在损伤已经不存在之后仍卡在开启状态——就像图~\ref{fig:bistable}中带有强反馈的神经元群。模型~\eqref{eq:sensitization}是一种简化；中枢敏化的存在已经测量。

\section{疼痛既是老师，也是中断}

除了报警，疼痛在机器里还有两项工作。

\emph{老师。}疼痛是负奖励：在误差公式~\eqref{eq:ctd}中，它以$r<0$的形式进入。第~\ref{sec:dofa}~章所讲的一切在这里同样适用：疼痛的先兆会提前引起误差，网络学会回避导致疼痛的事物。对人的实验表明，从疼痛中学习时，大脑的活动正是跟随时间差分误差~\cite{Seymour2004}，而且一部分\nt{DOPA}神经元也对不愉快的事件作出响应（第~\ref{sec:dofaalt}~章）。

\emph{中断。}疼痛把注意力从当前任务上拉开——这是它的本职，而不是副作用~\cite{EcclestonCrombez1999}。在第~\ref{sec:nora}~章中，我们曾就\nt{NORA}的短促簇放电讨论过同样的“中断”作用。而最快的响应甚至不到达大脑：从热物上缩手的动作直接在脊髓中闭合，所用的正是图~\ref{fig:antagonists}中的那对肌肉和对拮抗肌的同样否决。

\begin{keyblock}
\begin{proposition}[警报信号]
\label{prop:pain}
疼痛不是测量，而是高优先级的警报信号。它的适应远弱于测量通道，沿两根导线传送（快的报告“哪里”，慢的报告“有多糟”），通过一道由触摸关闭的闸门（而按闸门理论的假设，剧痛会打开它），其音量由广播通道从上方设定。除了疼痛打开闸门这一点，这些结构都已测量；简单模型~\eqref{eq:gate}和~\eqref{eq:sensitization}是简化。
\end{proposition}
\end{keyblock}

\section{小结}

\begin{table}[t]
\centering
\footnotesize
\setlength{\tabcolsep}{4pt}
\renewcommand{\arraystretch}{1.15}
\begin{tabular}{>{\raggedright}p{3.0cm}>{\raggedright}p{4.9cm}>{\raggedright\arraybackslash}p{4.9cm}}
\toprule
疼痛做什么 & 怎么做 & 已知程度 \\
\midrule
发出警报 & 高阈值，适应弱于触觉 & 阈值已测量；适应的比较是估计 \\
报告“哪里”和“有多糟” & 两根速度不同的导线~\eqref{eq:paintime} & 已测量 \\
通过闸门 & 经\nt{GABA}/\nt{GLYC}的否决~\eqref{eq:gate} & 闸门已测量；疼痛开门是理论假设；模型是简化 \\
改变音量 & 下行的\nt{SERO}、\nt{NORA}、阿片类 & 已测量；选择音量的规则是假说 \\
损伤后学习 & 增益增大~\eqref{eq:sensitization} & 敏化已测量；模型是简化 \\
教会回避 & 式~\eqref{eq:ctd}中的负奖励 & 与时间差分误差的相似性已测量 \\
中断工作 & 夺取注意力；缩回反射 & 已测量 \\
\bottomrule
\end{tabular}
\caption{疼痛作为警报信号。}
\label{tab:pain}
\end{table}

疼痛由我们已经见过的部件组成（表~\ref{tab:pain}）：\nt{GLUT}传感器、两根导线、经抑制性中间神经元的否决、广播式音量旋钮、慢RC电路、预测误差和中断。它只有一点是新的：它是机器唯一一个由机器自己设定音量的输入——一套主人既能调低、也能重新配置的警报系统。

\section{习题}

\begin{exercise}[\lvlA]
\label{ex:14:1}
\emph{两根导线。}信号来自脚部，$L=1$\,m。(a)~一次齐射沿一束同类导线传送，其速度布满范围~\eqref{eq:paintime}。对快、慢两类导线，齐射到达脊髓时被拉长多少？(b)~沿$15$和$1$\,m/s导线传来的两波疼痛，时间差达到$0.1$\,s才可分辨。距离多近时它们融为一体？
\end{exercise}

\begin{exercise}[\lvlB]
\label{ex:14:2}
\emph{触摸何时会痛。}闸门~\eqref{eq:gate}取正文的参数，增益$g$随敏化增大。$g$在多大以内，无论强度$\mu$多大，无痛的触摸都不会通过闸门？$g$为多少时，触摸$\mu=1$开始通过？
\end{exercise}

\begin{exercise}[\lvlB]
\label{ex:14:3}
\emph{敏化的稳定性。}输入恒定，闸门输出对$g$是线性的：$z=gA-C$，$A=\nu+w_{z\mu}\mu$，$C=\beta q\ge0$。(a)~求模型~\eqref{eq:sensitization}的平衡点$g^*$及其稳定条件。(b)~$k_s=0.5$时，分别在无触摸和触摸$\mu=1$时，求使模型失控的疼痛强度。
\end{exercise}

\begin{exercise}[\lvlA]
\label{ex:14:4}
\emph{疼痛的先兆。}时刻$0$的信号预示时刻$T=2$\,s的疼痛，在式~\eqref{eq:ctd}中$r(t)=-P\,\delta(t-T)$，$\tau_\gamma=2$\,s。(a)~求信号时刻$\delta$爆发的面积。(b)~在时刻$t_e=1$\,s的动作取消了疼痛。此刻$\delta$的爆发是多少？它会教给网络什么？
\end{exercise}

\begin{exercise}[\lvlB]
\label{ex:14:5}
\emph{自锁的警报。}给\texttt{models/\allowbreak pain\_gate.py}中的闸门加上动态~\eqref{eq:sensitization}和饱和$z\le4$。触摸$\mu=1$始终存在，损伤$\nu=2$持续时间$T$，然后$\nu=0$。用模拟求出$k_s=4$、$4.6$、$5$、$6$、$8$时，使警报保持开启的最小$T$（以$\tau_s$为单位）。
\end{exercise}

\begin{exercise}[\lvlB]
\label{ex:14:6}
\emph{设计闸门。}在$\beta=1.5$、$w_{z\mu}=0.3$、$g=1$时，选取式~\eqref{eq:gate}中的$q_0$、$w_{q\nu}$、$w_{q\mu}$，使无触摸时疼痛阈值为$0.2$，触摸$\mu=1$时为$1.0$，而在$\nu_\times=2.5$时触摸不再减轻疼痛。增益$g$在多大以内，无痛的触摸不会通过闸门？
\end{exercise}

\begin{exercise}[\lvlC]
\label{ex:14:7}
\emph{音量旋钮。}工作每单位时间带来价值$V$，带着损伤$\nu$工作的代价为$c\nu$，因此当$\nu>V/c$时中断是值得的；当$z>\theta=0.5$时疼痛引起中断。(a)~$V/c=0.25$、$0.5$、$2$时，无触摸的闸门~\eqref{eq:gate}取什么增益$g^*$才能给出这个阈值？(b)~机器可以根据本书中的哪些信号来估计$V$和$c$？
\end{exercise}
